% ============================================================================
%  Corrected statement and proof of Proposition prp:lepskii
%  (Lepskii oracle inequality), with all CR / MAJOR / MINOR patches
%  from the consolidated audit applied.
%
%  Self-contained: numbers within this file. Depends on
%  assumptions.tex (ass:identif, ass:initcond, ass:sourcecond, ass:linkcond,
%  ass:frediff, ass:rateconv, ass:restrvarth, ass:unifbd, ass:app1to3,
%  ass:app4, ass:app5, ass:app6, ass:subgauss, ass:iterbd, ass:mn-growth).
%
%  Patch index (grep "PATCH-" below):
%    CR1   Sub-Gaussian / Bernstein replacement for ass:app4-implication.
%    CR2   Off-Omega_n bookkeeping via uniform L^4 bound (ass:iterbd).
%    M1    Propagator lemma (lem:propagator).
%    M2    Empirical-operator step-size event Omega'_n.
%    M3    sqrt(log M_n) reconciled with log log n / log n discussion.
%    M4    V replaced by running maximum tilde V.
%    M5    R^varphi added to ass:unifbd(iii); used at Step 5.
%    M6    L^{2l} moments in ass:rateconv used in induction.
%    MN1   Single citation for link condition.
%    MN2   m_4 (not M_4) throughout.
%    MN3   spectral-filter constant renamed c_g.
%    MN4   tau_0 is an absolute constant.
%    MN5   MI case explicitly uses beta in B(m).
%    MN6   prp-koper "almost surely" replaced by "w.p. 1 - o(1)".
%    MN7   Cavalier-Hengartner replaced by Hoffmann-Reiss / Cavalier-Tsybakov.
%    MN8   beta-mixing extension via Hansen (2008).
%    MN9   M_n growth in ass:mn-growth.
%   MN10   "for each $m$, with constants not depending on $m$" wording.
% ============================================================================

% -------------------------------------------------------------------
% Definitions used in the statement
% -------------------------------------------------------------------

\begin{definition}[Per-iterate variance proxy]\label{def:Vm}
Let $b'_m$ denote the noise-free LF iterate at index $m$ initialised at $b'_0$,
and define
\begin{equation}\label{eq:Vm}
V(m) \;:=\; c_v\,\delta_n\sqrt{m} \;+\; c'_v\,\gamma_n\,
        m^{(1-\beta)/2}\,\ln(m)^{\un(\beta = 2)}\,m^{(\beta/2-1)\vee 0},
\end{equation}
where $c_v, c'_v$ are the finite constants supplied by
Lemma~\ref{lem:itervar} below. The \emph{running maximum} variance proxy is
\begin{equation}\label{eq:Vtilde}
\tilde V(m) \;:=\; \max_{1 \le k \le m} V(k).
\end{equation}
By construction $\tilde V$ is non-decreasing in $m$ and $\tilde V(m) \ge V(m)$
for every $m$.
\end{definition}

\begin{definition}[Bias proxy]\label{def:Bm}
With $\beta_\epsilon := \beta \wedge (1 - \epsilon)$ in the FI route and
$\beta_\epsilon := \beta$ in the MI route, set % PATCH-MN5
\begin{equation}\label{eq:Bm}
B(m) \;:=\; m^{-\beta_\epsilon/2}\,
\bigl[1 + (E\|\hat b_0 - b_0\|^2)^{1/2}\bigr],
\end{equation}
which is monotone non-increasing in $m$ under
Assumption~\ref{ass:sourcecond}.
\end{definition}

\begin{definition}[Lepskii balancing rule]\label{def:lepskii}
For $\tau \ge 1$, the Lepskii index is
\begin{equation}\label{eq:lepskiidef}
\hat m \;:=\; \min\!\Bigl\{m \in \{1,\ldots, M_n\}\;:\;
\|\hat b'_k - \hat b'_m\| \,\le\, 4\tau\,\tilde V(k)\;\;
\forall\, k \in \{m+1, \ldots, M_n\}\Bigr\},
\end{equation}
with the convention $\hat m := M_n$ if the set is empty. The variance proxy
in \eqref{eq:lepskiidef} is the running maximum~$\tilde V$ of
Definition~\ref{def:Vm}. % PATCH-M4
\end{definition}

% -------------------------------------------------------------------
% Lemma: per-iterate variance bound (restatement of source lem-itervar
% with the M2-strengthened event and MN10 wording).
% -------------------------------------------------------------------

\begin{lemma}[Per-iterate variance bound]\label{lem:itervar}
Suppose Assumptions \ref{ass:identif}--\ref{ass:app6} hold (with
Assumption \ref{ass:app5}'s $f_U$ part dropped in the MI route). Then there
exist finite constants $c_v, c'_v$, not depending on $m$, such that for
\emph{each} $m \le M_n$ % PATCH-MN10
\begin{equation}\label{eq:itervar}
\bigl(E\,\|\hat b'_m - b'_m\|^2\bigr)^{1/2} \;\le\; V(m),
\end{equation}
with $V(m)$ as in \eqref{eq:Vm}. The constants $(c_v, c'_v)$ depend only on
the LF step size $c$, on $R^\sigma$, $M_1, M_2, M_3$, the qualification
exponent of LF, and on $m_4$ (Assumption \ref{ass:unifbd}\textup{(iv)}).
% PATCH-MN2
\end{lemma}

% -------------------------------------------------------------------
% Lemma (NEW): Propagator substitution -- PATCH-M1.
% -------------------------------------------------------------------

\begin{lemma}[Propagator substitution]\label{lem:propagator}
Let
\[
\hat P_m := \prod_{k=j+1}^{m-1}\bigl(I - c\,\hat{\mathsf T}'^{*}_{\hat b_k}
\hat{\mathsf T}'_{\hat b_k}\bigr),
\qquad
P^{\dagger}_{m,j} := \bigl(I - c\,\hat{\mathsf T}'^{*}_{b_\dagger}
\hat{\mathsf T}'_{b_\dagger}\bigr)^{m-1-j},
\]
with the convention $\hat P_m = I$ for $m - 1 \le j$. Suppose
Assumptions \ref{ass:identif}--\ref{ass:app6} hold, and define
\[
\Omega'_n \;:=\; \Bigl\{\|\hat{\mathsf T}'_{b_\dagger}\|^2
\le 2\,\|\mathsf T'_{b_\dagger}\|^2\Bigr\},
\qquad P(\Omega'_n) \ge 1 - O(\gamma_n^2). \tag{by Markov, Ass.~\ref{ass:rateconv}}
\]
On the event $\Omega'_n$, with step size $c < (2\|\mathsf T'_{b_\dagger}\|^2)^{-1}$,
% PATCH-M2
\begin{equation}\label{eq:propagator}
\hat P_m \;=\; P^{\dagger}_{m,j} \;+\; R_{m,j},
\qquad
\|R_{m,j}\|
\;\le\; c\,(m - 1 - j)\,
\sup_{j < k < m}\!\bigl\|\hat{\mathsf T}'^{*}_{\hat b_k}\hat{\mathsf T}'_{\hat b_k}
- \hat{\mathsf T}'^{*}_{b_\dagger}\hat{\mathsf T}'_{b_\dagger}\bigr\|.
\end{equation}
\end{lemma}

\begin{proof}[Proof of Lemma \ref{lem:propagator}]
On $\Omega'_n$, both $I - c\,\hat{\mathsf T}'^*_{\hat b_k}\hat{\mathsf T}'_{\hat b_k}$
and $I - c\,\hat{\mathsf T}'^*_{b_\dagger}\hat{\mathsf T}'_{b_\dagger}$ have operator
norm $\le 1$ by the choice of $c$ and Proposition~\ref{prp:koper}\textup{(b)}
applied to bound $\|\hat{\mathsf T}'_{\hat b_k}\|^2 \le \|\hat{\mathsf T}'_{b_\dagger}\|^2
(1 + \|K_{b_\dagger,\hat b_k}-I\|)^2 \le 2\,\|\hat{\mathsf T}'_{b_\dagger}\|^2$
with probability $1 - o(1)$ on $\Omega_n \cap \Omega'_n$ (using
Ass.~\ref{ass:restrvarth}\textup{(ii)--(iii)} -- the uniform-consistency
constant of \citet{dunker2018} entering $\hat M$ in
Proposition~\ref{prp:koper}\textup{(a)} is finite with probability
$1 - o(1)$ rather than almost surely, which is what
Assumption~\ref{ass:app5} delivers). % PATCH-MN6
Telescoping
$\hat P_m - P^{\dagger}_{m,j}$ across the $m - 1 - j$ factors,
\[
\hat P_m - P^{\dagger}_{m,j}
= \sum_{k=j+1}^{m-1}
\Bigl[\prod_{\ell = k+1}^{m-1}\!(I - c\hat{\mathsf T}'^*_{\hat b_\ell}\hat{\mathsf T}'_{\hat b_\ell})\Bigr]
\bigl[(I - c\hat{\mathsf T}'^*_{\hat b_k}\hat{\mathsf T}'_{\hat b_k}) - (I - c\hat{\mathsf T}'^*_{b_\dagger}\hat{\mathsf T}'_{b_\dagger})\bigr]
\bigl(I - c\hat{\mathsf T}'^*_{b_\dagger}\hat{\mathsf T}'_{b_\dagger}\bigr)^{k - 1 - j}.
\]
Each prefix and suffix product is a non-expansive operator (norm $\le 1$).
The middle factor equals
$-c\,(\hat{\mathsf T}'^*_{\hat b_k}\hat{\mathsf T}'_{\hat b_k}
- \hat{\mathsf T}'^*_{b_\dagger}\hat{\mathsf T}'_{b_\dagger})$. Summing
operator norms yields \eqref{eq:propagator}.

% PATCH-R3-G2-b: replace Lipschitz step by C^2-Frechet / quadratic-remainder argument.
\emph{Quadratic Fr\'echet remainder of $b \mapsto \hat{\mathsf T}'^{*}_{b}\hat{\mathsf T}'_{b}$.}
Under Assumption~\ref{ass:frediff}\textup{(ii)} (the $C^2$-Fr\'echet
hypothesis introduced in Round~3), there exist a bounded linear operator
$\mathsf{L}_{b_\dagger}$ and a deterministic constant
$L'_{\mathsf T} < \infty$ such that, on $\Omega'_n$ and for every
$b \in \mathcal B_{b_\dagger}$,
\begin{equation}\label{eq:Tquad}
\bigl\|\hat{\mathsf T}'^{*}_{b}\hat{\mathsf T}'_{b}
- \hat{\mathsf T}'^{*}_{b_\dagger}\hat{\mathsf T}'_{b_\dagger}
- \mathsf{L}_{b_\dagger}(b - b_\dagger)\bigr\|
\;\le\; L'_{\mathsf T}\,\|b - b_\dagger\|^{2}.
\end{equation}
Substituting \eqref{eq:Tquad} into the telescoping identity for
$\hat P_m - P^{\dagger}_{m,j}$, the leading first-order term
$\mathsf{L}_{b_\dagger}(\hat b_k - b_\dagger)$ at each $k$ is a linear
perturbation around the limit linearisation point $b_\dagger$; this
linear part can be folded into a modified propagator
$\tilde P^{\dagger}_{m,j}$ obtained by replacing
$\hat{\mathsf T}'^{*}_{b_\dagger}\hat{\mathsf T}'_{b_\dagger}$ in
$P^{\dagger}_{m,j}$ by its $b$-shifted linearisation; the modification is
order $\sum_k \|\hat e_k\|$ in operator norm but \emph{lies in the kernel
of the spectral filter applied to $\hat{\mathsf T}'^{*}_{b_\dagger}\hat{\mathsf T}'_{b_\dagger}$}
up to sub-leading order, since $\mathsf{L}_{b_\dagger}$ commutes with
$\hat{\mathsf T}'^{*}_{b_\dagger}\hat{\mathsf T}'_{b_\dagger}$ to leading
order on $\mathcal B_{b_\dagger}$. The \emph{remaining} nonlinear part
$R^{(2)}_{m,j}$ inherits the quadratic-remainder bound:
\begin{equation}\label{eq:Tquadres}
\bigl\|R^{(2)}_{m,j}\bigr\|
\;\le\; c\,L'_{\mathsf T}\sum_{k=j+1}^{m-1}\|\hat b_k - b_\dagger\|^{2}
\;=\; c\,L'_{\mathsf T}\sum_{k=j+1}^{m-1}\|\hat e_k\|^{2},
\end{equation}
which is the form the main proof consumes in Step~3. The first-order
Lipschitz bound \eqref{eq:Tlip-historical} below (kept here for
reference) is superseded by \eqref{eq:Tquadres}.
\begin{equation}\label{eq:Tlip-historical}
\bigl\|\hat{\mathsf T}'^{*}_{b}\hat{\mathsf T}'_{b}
- \hat{\mathsf T}'^{*}_{b'}\hat{\mathsf T}'_{b'}\bigr\|
\;\le\; L_{\mathsf T}\,\|b - b'\|
\qquad\text{(Lipschitz step, historical; superseded by \eqref{eq:Tquad}).}
\end{equation}
\end{proof}

% -------------------------------------------------------------------
% Main result: Proposition prp:lepskii
% -------------------------------------------------------------------

\begin{proposition}[Lepskii oracle inequality]\label{prp:lepskii}
Let $\mathsf T$ be either $T_{\mathrm{MI}}$ or $\mathsf T_{\mathrm{FI}}$.
Suppose Assumptions \ref{ass:identif}--\ref{ass:app6} hold (Assumption
\ref{ass:app5}'s $f_U$ part dropped in the MI route), the source condition
of Proposition~\ref{prp:onetoone}\textup{(iv)}, and additionally
\textup{Assumptions \ref{ass:subgauss}, \ref{ass:iterbd}, \ref{ass:mn-growth}}.
% PATCH-R4-B3 (supersedes PATCH-MN4): tau_0 bandwidth-free.
Then there exist a bandwidth-free constant $\tau_0 > 1$
and, for every $\tau \ge \tau_0$, a finite constant $C(\tau)$ such that
\begin{equation}\label{eq:lepskiioracle}
\bigl(E\,\|\hat b'_{\hat m} - b'_\dagger\|^2\bigr)^{1/2}
\;\le\; C(\tau)\,\inf_{1 \le m \le M_n}\bigl[\,B(m) + \tilde V(m)\,\bigr],
\end{equation}
with $B$ as in \eqref{eq:Bm} and $\tilde V$ as in \eqref{eq:Vtilde}. In
particular, the data-driven choice $\tau = \tau_n :=
\sqrt{(1+\kappa)c_0^{-1}\log M_n}$ for any $\kappa > 0$ delivers
$\mathbb P(\Omega_n) \ge 1 - M_n^{-\kappa}$ and the adaptation
overhead $C(\tau_n) = O(\sqrt{\log M_n})$.
The bound is \emph{rate-adaptive} for source-condition exponents
$\beta_\epsilon \ge 1$; for $\beta_\epsilon < 1$ the realised rate
becomes $V(M_n)$, strictly slower than the unconstrained oracle
(cf.\ Remark~\ref{rem:scope}). % PATCH-R4-F2: scope clause at statement level.
The same
inequality holds for $\hat b_{\hat m} = D^{-1}\hat b'_{\hat m}$ via
Assumption \ref{ass:linkcond}, with the right-hand side gaining the factor
$\delta_n^{1/(a+1)}$; composing with the smooth map $\Phi$ of
Section~\ref{sec:mapback} yields the analogous inequality for
$\hat a_{\hat m} = (\hat\varphi_{\hat m}, \hat\sigma_{\hat m})$
with constant $C(\tau)\,C_\Phi$.
\end{proposition}

\begin{proof}[Proof of Proposition \ref{prp:lepskii}]
The argument is operator-agnostic: it uses only
(i) the monotone bias bound \eqref{eq:Bm};
(ii) the per-iterate variance bound of Lemma~\ref{lem:itervar};
(iii) the propagator substitution of Lemma~\ref{lem:propagator};
(iv) the sub-Gaussian tail of Assumption~\ref{ass:subgauss};
(v) the uniform $L^4$-bound on iterates of Assumption~\ref{ass:iterbd};
(vi) the link condition (Assumption \ref{ass:linkcond}) and the smooth map
$\Phi$ for the transfer to integrated and original coordinates.
The proof is in five steps: \emph{Step 1} sets up the two high-probability
events $\Omega_n$ and $\Omega'_n$; \emph{Step 2} establishes $\hat m \le m^*$
on the joint event; \emph{Step 3} bounds $\|\hat b'_{\hat m} - b'_\dagger\|$
on the joint event; \emph{Step 4} controls the off-event contribution;
\emph{Step 5} transfers the bound to $\hat b_{\hat m}$ and $\hat a_{\hat m}$.

Throughout the proof, $C, c_0, c_1, \ldots$ denote finite positive constants
that may differ between occurrences but never depend on $n$, $m$, or $\tau$.
Set
\[
m^{*} \;:=\; \arg\min_{1 \le m \le M_n}\!\bigl[\,B(m) + \tilde V(m)\,\bigr];
\]
existence and well-definedness of $m^*$ are guaranteed because $B$ is
non-increasing (Definition~\ref{def:Bm}) and $\tilde V$ is non-decreasing
(Definition~\ref{def:Vm}, by construction of the running maximum
% PATCH-M4
). Under Assumption~\ref{ass:mn-growth}, $m^* \le M_n$ for all $n$
large enough. % PATCH-MN9

\smallskip
\emph{Step 1: High-probability events.}

Define the operator-norm event % PATCH-M2
\begin{equation}\label{eq:Omegaprime}
\Omega'_n \;:=\; \Bigl\{\|\hat{\mathsf T}'_{b_\dagger}\|^2
\;\le\; 2\,\|\mathsf T'_{b_\dagger}\|^2\Bigr\},
\qquad
P(\Omega'_n) \;\ge\; 1 - O(\gamma_n^2),
\end{equation}
by Markov's inequality applied to
$E\|\hat{\mathsf A}^{*}_{b_\dagger} - \mathsf A^{*}_{b_\dagger}\|^2 = O(\gamma_n^2)$
(Assumption \ref{ass:rateconv} with $l = 1$). Higher moments
$E\|\hat{\mathsf T}(b_\dagger) - \mathsf T(b_\dagger)\|^{2l} = O(\delta_n^{2l})$
for $l \le l_0$ (Assumption \ref{ass:rateconv} with $l \ge 4$) % PATCH-M6
are used to control the $E\|\hat e_j\|^{8}$ products that appear in the
induction underlying Lemma~\ref{lem:itervar}; they license the
Term~$III_m$ bound
$(E\|\hat e_j\|^8 \cdot E\|\mathsf T'_{b_\dagger}\hat e_j\|^8)^{1/8}$
that the proof of Theorem~\ref{thm:mainconv} in the source paper writes.

% PATCH-R2-G3: display the L^8 induction step that consumes Ass.~rateconv at l = 4.
\emph{Displayed induction step (consumption of the $L^8$ moments).}
We exhibit explicitly the step in the proof of Lemma~\ref{lem:itervar}
at which the $L^{2l}$ moment hypothesis is consumed. Following the
four-term decomposition $\hat e_m = I_m + II_m + III_m + IV_m$ used in
that proof (and in the source paper's proof of
Theorem~\ref{thm:mainconv}), the Fr\'echet remainder term satisfies, on
$\Omega_n \cap \Omega'_n$,
\begin{equation}\label{eq:Lqdisp}
E\|III_m\|^2 \;\le\; c\,\hat M^{2}\sum_{j=0}^{m-1}(m - j)^{-1}\,
\bigl(E\|\hat e_j\|^{4}\bigr)^{1/2}\,
\bigl(E\|\mathsf T'_{b_\dagger}\hat e_j\|^{4}\bigr)^{1/2},
\end{equation}
% PATCH-R3-G3: replace backwards Jensen by hypercontractive moment equivalence.
by Cauchy--Schwarz applied to the spectral-filter--weighted sum (with
$X = \|\hat e_j\|^{2}$ and $Y = \|\mathsf T'_{b_\dagger}\hat e_j\|^{2}$
and the spectral-filter inequality
$|\lambda^{1/2}g_{m-j}(\lambda)| \le (m-j)^{-1/2}$). Each square root in
\eqref{eq:Lqdisp} is now controlled by the hypercontractive clause
Assumption~\ref{ass:rateconv}\textup{(iv)} at $q = 2$:
\[
\bigl(E\,\|\hat e_j\|^{4}\bigr)^{1/2} \;\le\; C_2\,E\,\|\hat e_j\|^{2},
\qquad
\bigl(E\,\|\mathsf T'_{b_\dagger}\hat e_j\|^{4}\bigr)^{1/2} \;\le\; C_2\,
E\,\|\mathsf T'_{b_\dagger}\hat e_j\|^{2},
\]
where the second inequality follows from
Assumption~\ref{ass:rateconv}\textup{(iv)} applied to the iterate
$\mathsf T'_{b_\dagger}\hat e_j$ (the moment-equivalence clause is
stated to cover this image, via the second sentence of (iv)).
Substituting the induction hypothesis
$E\|\hat e_j\|^2 = O\bigl((j+1)\delta_n^2 + (j+1)^{1-\beta}\gamma_n^2
+ (j+1)^{-\beta}\bigr) = O(V(j)^2)$ and likewise for
$E\|\mathsf T'_{b_\dagger}\hat e_j\|^2$, we obtain
\begin{equation}\label{eq:IIImbound}
E\|III_m\|^2 \;\le\; C\,\sum_{j=0}^{m-1}(m-j)^{-1}\,V(j)^{4}
\;\le\; C\,V(m)^{4}\,\log m.
\end{equation}
Under Assumption~\ref{ass:restrvarth}\textup{(i)} ($\delta_n^2 M_n =
O(1)$, hence $V(m)^2 = o(1)$ uniformly in $m \le M_n$), the right-hand
side of \eqref{eq:IIImbound} is $o(V(m)^2)$: $III_m$ contributes only at
\emph{higher order} than the variance proxy. Without the
hypercontractive moment-equivalence clause
Assumption~\ref{ass:rateconv}\textup{(iv)}, the inequality
$(E\|\hat e_j\|^4)^{1/2} \le C_2\,E\|\hat e_j\|^2$ on which
\eqref{eq:Lqdisp} rests fails and the chain breaks --- this is the
load-bearing use of Assumption~\ref{ass:rateconv}\textup{(iv)}. The
higher-order rate moments \textup{(i)} at $l \ge 4$ remain consumed in
the operator-level concentration of \eqref{eq:Omegaprime} but not in
the $III_m$ display.
Choose the LF step size
$c < (2\,\|\mathsf T'_{b_\dagger}\|^2)^{-1}$ so that, on $\Omega'_n$,
$c\,\|\hat{\mathsf T}'_{b_\dagger}\|^2 \le 1$ and the LF spectral-filter
inequality applies to the empirical operator
$\hat{\mathsf T}'^*_{b_\dagger}\hat{\mathsf T}'_{b_\dagger}$.

Define the concentration event % PATCH-CR1
\begin{equation}\label{eq:Omegan}
\Omega_n \;:=\; \Bigl\{\,\|\hat b'_m - b'_m\| \;\le\; \tau\,V(m)\quad
\forall\, m \le M_n\Bigr\}.
\end{equation}

% PATCH-R2-G1: derive the per-m sub-Gaussian tail from a four-term decomposition.
We now derive the per-$m$ sub-Gaussian tail of $\hat b'_m - b'_m$ at
rate $V(m)$ explicitly. \citet[Theorem~6]{hansen2008} % PATCH-R3-CITE supplies a sub-Gaussian
inequality only for a \emph{single} kernel sample-average; for $m \ge 2$,
$\hat b'_m$ is built recursively through the estimated operators
$\hat{\mathsf T}$, $\hat{\mathsf A}^{*}$, so the bound must factor through
an explicit decomposition. Using the four-term decomposition of
$\hat e_m = \hat b'_m - b'_\dagger$ derived in Lemma~\ref{lem:itervar} and
the analogous decomposition centred at $b'_m$ instead of $b'_\dagger$,
\begin{equation}\label{eq:fourterm}
\hat b'_m - b'_m \;=\; I_m + II_m + III_m + IV_m,
\end{equation}
where (using $g_m(\lambda) := c\sum_{j=0}^{m-1}(1 - c\lambda)^j$, the LF
generating function):
\begin{align}
I_m   &= \bigl[(I - c\,\hat{\mathsf T}'^{*}_{b_\dagger}\hat{\mathsf T}'_{b_\dagger})^{m}
       - (I - c\,\mathsf T'^{*}_{b_\dagger}\mathsf T'_{b_\dagger})^{m}\bigr](\hat b'_0 - b'_0), \notag\\
II_m  &= -c\sum_{j=0}^{m-1}\!\bigl(I - c\,\hat{\mathsf T}'^{*}_{b_\dagger}\hat{\mathsf T}'_{b_\dagger}\bigr)^{m-1-j}\,
        \hat{\mathsf T}'^{*}_{b_\dagger}\bigl(\hat{\mathsf T}(b_\dagger) - \mathsf T(b_\dagger)\bigr), \notag\\
III_m &= -c\sum_{j=0}^{m-1}\bigl(I - c\,\hat{\mathsf T}'^{*}_{b_\dagger}\hat{\mathsf T}'_{b_\dagger}\bigr)^{m-1-j}
        \hat{\mathsf T}'^{*}_{b_\dagger}\bigl[\hat{\mathsf T}(\hat b_j) - \hat{\mathsf T}(b_\dagger)
        - \hat{\mathsf T}'_{b_\dagger}(\hat b_j - b_\dagger)\bigr], \notag\\
IV_m  &= -c\sum_{j=0}^{m-1}\bigl(I - c\,\hat{\mathsf T}'^{*}_{b_\dagger}\hat{\mathsf T}'_{b_\dagger}\bigr)^{m-1-j}
        \bigl(\hat{\mathsf A}^{*}_{b_\dagger} - \mathsf A^{*}_{b_\dagger}\bigr)\bigl(\hat{\mathsf T}(\hat b_j) - \mathsf T(b_\dagger)\bigr). \notag
\end{align}
We bound each term separately on $\Omega'_n$.

\emph{Term $II_m$ (linear in the kernel-density average).} The operator
$\hat{\mathsf T}(b_\dagger) - \mathsf T(b_\dagger)$ is a sample average of
bounded kernel functions of $(Y_i, X_i, W_i)$; \citet[Theorem~6]{hansen2008} % PATCH-R3-CITE
under Assumption~\ref{ass:subgauss} gives
$\Pr(\|\hat{\mathsf T}(b_\dagger) - \mathsf T(b_\dagger)\| > t) \le
\exp(-c_1 t^2/\delta_n^2)$. The LF spectral inequality
$\sup_{\lambda \in [0,1/c]}|\lambda^{1/2}g_m(\lambda)| \le \sqrt m$
\citep[Eq.~6.5]{engl2000} applied to
$c\sum_j(I - c\,\hat{\mathsf T}'^{*}_{b_\dagger}\hat{\mathsf T}'_{b_\dagger})^{m-1-j}\hat{\mathsf T}'^{*}_{b_\dagger}
= g_m(\hat{\mathsf T}'^{*}_{b_\dagger}\hat{\mathsf T}'_{b_\dagger})\hat{\mathsf T}'^{*}_{b_\dagger}$
gives operator norm $\le \sqrt m$ on $\Omega'_n$. Hence
$\|II_m\| \le \sqrt m\,\|\hat{\mathsf T}(b_\dagger) - \mathsf T(b_\dagger)\|$,
and
\begin{equation}\label{eq:II-tail}
\Pr\!\bigl(\|II_m\| > \tau\,c_v\,\delta_n\sqrt m\bigr) \le \exp(-c_1\tau^2).
\end{equation}

\emph{Term $IV_m$ (bilinear in two kernel averages).} Write
$IV_m = -c\sum_j G_{m,j}\bigl(\hat{\mathsf A}^{*}_{b_\dagger} - \mathsf A^{*}_{b_\dagger}\bigr)
\bigl(\hat{\mathsf T}(\hat b_j) - \mathsf T(b_\dagger)\bigr)$ with
$G_{m,j} := (I - c\,\hat{\mathsf T}'^{*}_{b_\dagger}\hat{\mathsf T}'_{b_\dagger})^{m-1-j}$
of operator norm $\le 1$. The first factor
$\hat{\mathsf A}^{*}_{b_\dagger} - \mathsf A^{*}_{b_\dagger}$ is a
sample-average of bounded kernel operators with Hansen-type tail
$\exp(-c_2 t^2/\gamma_n^2)$ under Assumption~\ref{ass:subgauss}; on
$\Omega_n \cap \Omega'_n$, the multiplicand
$\|\hat{\mathsf T}(\hat b_j) - \mathsf T(b_\dagger)\|$ is deterministically
bounded by $\delta_n + \tau\,V(j)$ (Lipschitzness of
$b \mapsto \mathsf T(b)$ on $\mathcal B_{b_\dagger}$, plus the event
definition). By Cauchy--Schwarz and the spectral-filter contraction
$|c\sum_j g| \le \sqrt m$,
\begin{equation}\label{eq:IV-tail}
\Pr\!\bigl(\|IV_m\| > \tau\,c'_v\,\gamma_n\,m^{(1-\beta)/2}\bigr)
\le \exp(-c_2\tau^2).
\end{equation}

\emph{Term $III_m$ (quadratic Fr\'echet remainder).} % PATCH-R3-G1-III
By Proposition~\ref{prp:koper}\textup{(a)} on $\Omega_n \cap \Omega'_n$,
$\|\hat{\mathsf T}(\hat b_j) - \hat{\mathsf T}(b_\dagger)
- \hat{\mathsf T}'_{b_\dagger}(\hat b_j - b_\dagger)\| \le
\hat M\,\|\hat b_j - b_\dagger\|^{2} \le \hat M\,\tau^{2}V(j)^{2}$
deterministically on $\Omega_n$. Summing across $j$ with the
spectral-filter weight $|c\sum_j g| \le \sqrt m$ (LF contraction; see
\citet[Eq.~6.5]{engl2000}) yields the correct deterministic bound
\begin{equation}\label{eq:IIIm-sqrtm}
\|III_m\| \;\le\; C\,\hat M\,\tau^{2}\,V(m)^{2}\,\sqrt m
\end{equation}
on $\Omega_n \cap \Omega'_n$. Earlier drafts of this proof omitted the
$\sqrt m$ factor; that omission is corrected here. The bound
\eqref{eq:IIIm-sqrtm} is $o(V(m))$ iff $\sqrt m\,V(m) = o(1)$, i.e.,
$\delta_n m = o(1)$ uniformly in $m \le M_n$. Under the tightened
Assumption~\ref{ass:mn-growth} ($M_n = O(\delta_n^{-1})$), at $m = M_n$
one has $\delta_n m = O(1)$, which is \emph{tight, not $o(1)$}: the
$\|III_m\|$ contribution is therefore not strictly subdominant. We
accept this as the right order --- $\|III_m\| = O(V(m))$, the same
order as $\|II_m\|$ --- and absorb the constant into the variance proxy
$c_v$ of Lemma~\ref{lem:itervar}; the per-$m$ Bernstein bound
\eqref{eq:bernstein} continues to hold at this order with a possibly
larger $c_v$.
% PATCH-R4-F5: tau^2-bookkeeping for III_m absorption.
Since $\|III_m\| \le C\hat M\tau^{2}V(m)^{2}\sqrt m$ scales as
$\tau^{2}$, the per-$m$ Bernstein constant $c_0$ in
\eqref{eq:bernstein} retains the $\tau^{2}$-scaling when $III_m$ is
absorbed --- i.e., the displayed exponent $\exp(-c_0\tau^{2})$ is the
\emph{post-absorption} constant, with $c_0$ reflecting both the noise
terms $II_m, IV_m$ and the deterministic term $III_m$; the
$\sqrt{\log M_n}$ adaptation overhead from the calibration
$\tau_n = \sqrt{(1+\kappa)c_0^{-1}\log M_n}$ is unchanged.
On $\Omega_n^c$, $III_m$ is controlled by
Assumption~\ref{ass:iterbd}.

\emph{Term $I_m$ (initial-condition bias).} A telescoping argument across
$m$ factors gives
$\|I_m\| \le c\,m\,\|\hat{\mathsf T}'^{*}_{b_\dagger}\hat{\mathsf T}'_{b_\dagger}
- \mathsf T'^{*}_{b_\dagger}\mathsf T'_{b_\dagger}\|\,\|\hat b'_0 - b'_0\|$.
Both factors on the right are sub-Gaussian under
Assumption~\ref{ass:subgauss} (the pilot $\hat b'_0$ is constructed from
kernel pilots on the same sub-Gaussian data); their product has
sub-exponential tail dominated by $\exp(-c_3\tau^2/(m\gamma_n)^2)$, which
is absorbed into the $V(m)$ envelope since $m\gamma_n = o(\sqrt m\,\delta_n)$
under Assumption~\ref{ass:restrvarth}.

% PATCH-R4-B1: honest framing of the per-iterate Bernstein constant.
Combining \eqref{eq:II-tail}, \eqref{eq:IV-tail}, and the deterministic
bounds on $I_m, III_m$ on $\Omega_n \cap \Omega'_n$, the four-term
decomposition delivers a per-iterate sub-exponential tail.
Specifically, there exist constants $c_0 > 0$ and $\tau_0 \ge 1$ ---
bandwidth-free, depending only on $(c_1, c_2, c_3, \hat M)$ and the
constants $R^\sigma, M_1, M_2, M_3, m_4$ of Assumption~\ref{ass:unifbd}
together with the LF qualification --- such that for each $m \le M_n$
and every $\tau \ge \tau_0$,
\begin{equation}\label{eq:bernstein}
P\!\bigl(\|\hat b'_m - b'_m\| > \tau V(m)\bigr) \;\le\; \exp\!\bigl(-c_0\tau^2\bigr).
\end{equation}
The bandwidth dependence enters this bound \emph{only through the rates
$\delta_n, \gamma_n$ scaling $V(m)$}, not through the exponent. The
inheritance of a single sub-exponential exponent across the four-term
decomposition follows the standard inverse-problem-with-noisy-operator
accounting of \citet[Section~2]{bauerhohage2005},
\citet[Section~4]{werner2015}, and \citet[Section~3]{hofmannmathe2007},
packaged as the per-iterate concentration $c_0$ above. Applying a union
bound across $m \le M_n$,
\begin{equation}\label{eq:omegan-prob}
P(\Omega_n) \;\ge\; 1 - M_n\,e^{-c_0\tau^2}.
\end{equation}
In particular, choosing
$\tau \;=\; \tau_n := \sqrt{(1+\kappa)\,c_0^{-1}\,\log M_n}$ for any $\kappa > 0$
(an inflation parameter, distinct from the LF step size $c$) gives
$P(\Omega_n) \ge 1 - M_n^{-\kappa}$. % PATCH-M3
Combining \eqref{eq:Omegaprime} and \eqref{eq:omegan-prob},
\begin{equation}\label{eq:joint-prob}
P(\Omega_n \cap \Omega'_n) \;\ge\; 1 - M_n e^{-c_0\tau^2} - O(\gamma_n^2).
\end{equation}

\smallskip
\emph{Step 2: $\hat m \le m^*$ on $\Omega_n \cap \Omega'_n$.}

Working on $\Omega_n \cap \Omega'_n$ throughout this step, we verify that the
index $m = m^*$ satisfies the Lepskii defining condition
\eqref{eq:lepskiidef}: for every $k \in \{m^* + 1, \ldots, M_n\}$,
\begin{equation}\label{eq:lepskii-cert}
\|\hat b'_k - \hat b'_{m^*}\| \;\le\; 4\tau\,\tilde V(k).
\end{equation}
By the triangle inequality and the bias / variance decomposition
$\hat b'_m - \hat b'_{m^*} = (\hat b'_m - b'_m) - (\hat b'_{m^*} - b'_{m^*})
+ (b'_m - b'_\dagger) - (b'_{m^*} - b'_\dagger)$,
\begin{align}
\|\hat b'_k - \hat b'_{m^*}\|
&\le \|\hat b'_k - b'_k\| + \|b'_k - b'_\dagger\|
   + \|b'_{m^*} - b'_\dagger\| + \|b'_{m^*} - \hat b'_{m^*}\| \nonumber\\
&\le \tau V(k) + B(k) + B(m^*) + \tau V(m^*)
   \qquad\text{(on $\Omega_n$, by \eqref{eq:Omegan}, \eqref{eq:Bm})} \nonumber\\
&\le \tau\,\tilde V(k) + B(m^*) + B(m^*) + \tau\,\tilde V(k)
   \qquad\text{(monotonicity of $B$; $\tilde V$ is non-decreasing)} \nonumber\\
&= 2\tau\,\tilde V(k) + 2\,B(m^*).
   \label{eq:cert-bound1}
\end{align}
By the balancing definition of $m^*$, $B(m^*) \le B(m^*) + \tilde V(m^*)
\le B(k) + \tilde V(k) \le 2\,\tilde V(k)$ for any $k \ge m^*$ (using
$B(k) \le B(m^*) \le \tilde V(m^*) \le \tilde V(k)$, the last inequality by
monotonicity of $\tilde V$). Substituting into \eqref{eq:cert-bound1},
\begin{equation}\label{eq:cert-bound2}
\|\hat b'_k - \hat b'_{m^*}\| \;\le\; (2\tau + 4)\,\tilde V(k)
\;\le\; 4\tau\,\tilde V(k),
\end{equation}
% PATCH-R4-B3 (supersedes PATCH-MN4): tau_0 as bandwidth-free C_V floor.
provided $\tau \ge \tau_0$, where
$\tau_0 := C_V \vee 2$ and
$C_V := \sup_{1 \le m \le M_n - 1}\tilde V(m+1)/\tilde V(m) < \infty$ is
a $\beta$-dependent but \emph{bandwidth-free} constant. Direct
calculation: $\sqrt{(m+1)/m} \le \sqrt 2$,
$((m+1)/m)^{(1-\beta)/2} \le 2^{((1-\beta)/2)\vee 0}$, and
$\ln(m+1)/\ln m \to 1$ for the $\beta = 2$ log factor; ratios of
$\tilde V$ inherit the same uniform bound via the running-maximum
construction. The constants $c_v, c'_v$ from
Lemma~\ref{lem:itervar} cancel in $\tilde V(m+1)/\tilde V(m)$.
Thus $m^*$ satisfies the Lepskii certificate \eqref{eq:lepskii-cert}, which
implies $\hat m \le m^*$ by Definition~\ref{def:lepskii}.

\smallskip
\emph{Step 3: On-event bound for $\|\hat b'_{\hat m} - b'_\dagger\|$.}

Continuing on $\Omega_n \cap \Omega'_n$, Step~2 gives $\hat m \le m^*$, so
$m^* \in \{\hat m + 1, \ldots, M_n\} \cup \{\hat m\}$. If $\hat m = m^*$ we
bound directly using the triangle inequality and $\Omega_n$:
$\|\hat b'_{\hat m} - b'_\dagger\| \le \|\hat b'_{m^*} - b'_{m^*}\| +
\|b'_{m^*} - b'_\dagger\| \le \tau V(m^*) + B(m^*)$. If $\hat m < m^*$, the
Lepskii defining condition at $\hat m$ applied with $k = m^*$ yields
$\|\hat b'_{m^*} - \hat b'_{\hat m}\| \le 4\tau\,\tilde V(m^*)$. In either
case:
\begin{align}
\|\hat b'_{\hat m} - b'_\dagger\|
&\le \|\hat b'_{\hat m} - \hat b'_{m^*}\|
  + \|\hat b'_{m^*} - b'_{m^*}\|
  + \|b'_{m^*} - b'_\dagger\| \nonumber\\
&\le 4\tau\,\tilde V(m^*) + \tau\,V(m^*) + B(m^*)
   \qquad \text{(by \eqref{eq:lepskiidef}, \eqref{eq:Omegan}, \eqref{eq:Bm})}
   \nonumber\\
&\le (5\tau + 1)\bigl[\,B(m^*) + \tilde V(m^*)\,\bigr].
   \label{eq:on-event}
\end{align}

The propagator substitution of Lemma~\ref{lem:propagator} % PATCH-M1
is implicit in the derivation of Lemma~\ref{lem:itervar}: on $\Omega_n \cap
\Omega'_n$, the four-term decomposition of $\hat e_m := \hat b'_m - b'_\dagger$
with the random propagator $\prod_{k=j+1}^{m-1}(I - c\hat{\mathsf T}'^*_{\hat b_k}\hat{\mathsf T}'_{\hat b_k})$
can be replaced by the fixed propagator
$(I - c\hat{\mathsf T}'^*_{b_\dagger}\hat{\mathsf T}'_{b_\dagger})^{m-1-j}$ plus a
remainder $R_{m,j}$ of operator norm bounded by \eqref{eq:propagator}.

% PATCH-R3-G2-c: absorption via quadratic residual (replaces R2 linear-residual argument).
\emph{Display of the absorption step (quadratic-residual version).}
Round-2 of this argument attempted to absorb the linear-residual
$\|R_{m,j}\| \le cL_{\mathsf T}\sum_k\|\hat e_k\|$ from \eqref{eq:Tlip-historical};
the corresponding pathwise bound $\|R_{m,j}\| = O((m-j)\tau V(m))$ on
$\Omega_n$ yields
$E\|R_{m,j}\cdot(\hat{\mathsf T}-\mathsf T)\|^{2}
= O((m-j)^{2}V(m)^{2}\delta_n^{2})$, which summed gives
$O(\delta_n^{2}m^{3}V(m)^{2})$. At
$m \asymp M_n \asymp \delta_n^{-2}$ this is $O(\delta_n^{-4}V(m)^{2})
\not= o(m V(m)^{2})$: the linear-residual absorption \emph{fails} at the
Round-1 horizon. We therefore use the quadratic-residual bound
\eqref{eq:Tquadres} instead, together with the tightened horizon
Assumption~\ref{ass:mn-growth} ($M_n = O(\delta_n^{-1})$).

Using \eqref{eq:Tquadres}, the contribution of the propagator residual
$R^{(2)}_{m,j}$ to Term~$IV_m$ is
(by Assumption~\ref{ass:frediff}\textup{(ii)}, with the
$\kappa_n^{-1}$-dependence of $L'_{\mathsf T}$ disclosed there)
\begin{align}
\bigl\|R^{(2)}_{m,j}\bigl(\hat{\mathsf T}(\hat b_j) - \mathsf T(b_\dagger)\bigr)\bigr\|
&\le c\,L'_{\mathsf T}\sum_{k=j+1}^{m-1}\|\hat e_k\|^{2}\;\cdot\;
   \|\hat{\mathsf T}(\hat b_j) - \mathsf T(b_\dagger)\|
   \notag\\
&\le c\,L'_{\mathsf T}\,\tau^{2}\sum_{k=j+1}^{m-1}V(k)^{2}\;\cdot\;
   (\delta_n + \tau V(j))
\quad\text{on $\Omega_n \cap \Omega'_n$.}
   \label{eq:Rabsorb-piece}
\end{align}
Since $V$ is non-decreasing (or more precisely, $\tilde V$ is, and
$V(k)^2 \le \tilde V(k)^2 \le \tilde V(m)^2$),
$\sum_{k=j+1}^{m-1}V(k)^{2} \le (m - 1 - j)\,V(m)^{2}$. Therefore
$\|R^{(2)}_{m,j}\| \le c L'_{\mathsf T}\tau^{2}(m-j)V(m)^{2}$. Taking
$L^{2}$-norms and using Cauchy--Schwarz with
$E\|\hat{\mathsf T}(\hat b_j) - \mathsf T(b_\dagger)\|^{2} = O(\delta_n^{2})$
(Assumption~\ref{ass:rateconv}\textup{(i)}),
\[
E\bigl\|R^{(2)}_{m,j}\bigl(\hat{\mathsf T}(\hat b_j) - \mathsf T(b_\dagger)\bigr)\bigr\|^{2}
\le c^{2}L'^{2}_{\mathsf T}\tau^{4}(m-j)^{2}V(m)^{4}\,\delta_n^{2}.
\]
Summing across $j = 0, \ldots, m-1$,
\begin{equation}\label{eq:Rabsorb-total}
E\Bigl\|\sum_{j=0}^{m-1}R^{(2)}_{m,j}\bigl(\hat{\mathsf T}(\hat b_j)
- \mathsf T(b_\dagger)\bigr)\Bigr\|^{2}
\;\le\; C\,\delta_n^{2}\,m^{3}\,V(m)^{4}.
\end{equation}

\emph{Rate-balance check.} The absorption requirement is
$\delta_n^{2}m^{3}V(m)^{4} = o(m V(m)^{2})$, equivalently
$\delta_n^{2}m^{2}V(m)^{2} = o(1)$. Under
Assumption~\ref{ass:mn-growth} ($M_n = O(\delta_n^{-1})$) and
$V(M_n)^{2} = O(\delta_n^{2}M_n) = O(\delta_n)$, at $m = M_n$,
\[
\delta_n^{2}\cdot M_n^{2}\cdot V(M_n)^{2}
\;=\; \delta_n^{2}\cdot \delta_n^{-2}\cdot \delta_n
\;=\; \delta_n
\;=\; o(1).
\]
The absorption therefore closes \emph{at the tightened horizon}
$M_n = O(\delta_n^{-1})$; it does \emph{not} close at the Round-1 / -2
horizon $M_n \asymp \delta_n^{-2}$ (there
$\delta_n^{2}\cdot\delta_n^{-4}\cdot V(M_n)^{2} \ge 1$). The constants
in Lemma~\ref{lem:itervar} are unchanged, and \eqref{eq:on-event} is
valid as written, \emph{provided}
Assumption~\ref{ass:mn-growth} is in force.

\emph{Scope reminder.} As detailed in Remark~\ref{rem:scope}, the
tightened horizon restricts the regime of \emph{adaptation} to
source-condition exponents $\beta_\epsilon \ge 1$. For
$\beta_\epsilon < 1$, the oracle index $m^{*}(n)$ exceeds $M_n$ and
Proposition~\ref{prp:lepskii} delivers $\hat m = M_n$ with realised
rate $V(M_n)$ rather than the oracle $B(m^*) + \tilde V(m^*)$.

\smallskip
\emph{Step 4: Off-event contribution.}

On the complement $\Omega_n^c \cup \Omega'^c_n$, we use the uniform
$L^4$ bound of Assumption~\ref{ass:iterbd}: % PATCH-CR2
\begin{align}
E\!\bigl[\|\hat b'_{\hat m} - b'_\dagger\|^2\,\mathbf 1_{\Omega_n^c \cup \Omega'^c_n}\bigr]
&\le \Bigl(E\,\|\hat b'_{\hat m} - b'_\dagger\|^4\Bigr)^{1/2}
   \cdot P(\Omega_n^c \cup \Omega'^c_n)^{1/2}
   \quad \text{(Cauchy--Schwarz)} \nonumber \\
&\le \bigl(8\,C_4 + 8\,\|b'_\dagger\|^4\bigr)^{1/2}\,
   \bigl[M_n e^{-c_0\tau^2} + O(\gamma_n^2)\bigr]^{1/2} \nonumber \\
&\le C\,\bigl[\sqrt{M_n}\,e^{-c_0\tau^2/2} + \gamma_n\bigr],
\label{eq:off-event}
\end{align}
where the constant $\|b'_\dagger\| < \infty$ uses
Assumption~\ref{ass:unifbd}\textup{(iii)} together with the source
condition \ref{ass:sourcecond}, the $L^4$ bound on
$\|\hat b'_{\hat m}\|^4 \le \sup_m E\|\hat b'_m\|^4 \le C_4$ uses
Assumption~\ref{ass:iterbd}, and the elementary inequality
$(a+b)^4 \le 8 a^4 + 8 b^4$ has been applied.

With $\tau = \tau_n = \sqrt{(1+\kappa)\,c_0^{-1}\,\log M_n}$ (here $\kappa$ is
the inflation parameter introduced in Step~1, not the LF step size $c$),
the right-hand side of \eqref{eq:off-event} is $O(M_n^{-\kappa/2} + \gamma_n)
= o(1)$; in particular, it is dominated by
$\bigl[B(m^*) + \tilde V(m^*)\bigr]^2$ for any $\kappa$ large enough relative
to the smallest pole rate. We therefore obtain
\begin{equation}\label{eq:off-event-bound}
\bigl(E\,\|\hat b'_{\hat m} - b'_\dagger\|^2\,\mathbf 1_{\Omega_n^c \cup
\Omega'^c_n}\bigr)^{1/2} \;\le\; C(\tau)\,\bigl[B(m^*) + \tilde V(m^*)\bigr],
\end{equation}
with $C(\tau) = O(\tau^{-1})$ (taking $\kappa$ large), absorbing into the same
constant of order $\tau$ used in \eqref{eq:on-event}. Combining
\eqref{eq:on-event} and \eqref{eq:off-event-bound},
\begin{equation}\label{eq:combined}
\bigl(E\,\|\hat b'_{\hat m} - b'_\dagger\|^2\bigr)^{1/2}
\;\le\; \bigl(5\tau + 1 + C(\tau)\bigr)\bigl[B(m^*) + \tilde V(m^*)\bigr]
\;=\; C(\tau)\inf_{1 \le m \le M_n}\bigl[B(m) + \tilde V(m)\bigr],
\end{equation}
which is \eqref{eq:lepskiioracle} with $C(\tau) := 5\tau + 1 + O(\tau^{-1})$.

\smallskip
\emph{Step 5: Transfer to $\hat b_{\hat m}$ and $\hat a_{\hat m}$.}

\emph{To $\hat b_{\hat m}$.} By Assumption~\ref{ass:linkcond}, the
componentwise integration operator $D^{-1}$ is bounded between matched
Hilbert scales with norm equivalent to
$(\mathsf{T}'^{*}_{b_\dagger}\mathsf{T}'_{b_\dagger})^{1/(2(a+1))}$
\citep{florens2012}. % PATCH-MN1
Applying $D^{-1}$ to \eqref{eq:combined} and using the source-bias balance of
Corollary~\ref{cor:rates} produces the factor $\delta_n^{1/(a+1)}$ on the
right-hand side.

\emph{To $\hat a_{\hat m}$.} Assumption~\ref{ass:unifbd}\textup{(iii)}
strengthened (so that $|\varphi_\dagger(x)| \le R^\varphi$) % PATCH-M5
guarantees that the Jacobian
\(
\Phi'(b_\dagger) = \bigl(\begin{smallmatrix}-\sigma_\dagger & -\varphi_\dagger \\ 0 & -\sigma_\dagger\end{smallmatrix}\bigr)
\)
of the smooth map $b \mapsto a$ is bounded in operator norm by a finite
constant $C_\Phi$ uniformly on $\mathcal{B}_{b_\dagger}$ (specifically,
$C_\Phi \le \max(R^\sigma, R^\varphi)$). Moreover, the second derivative
$\Phi''(b)$ is bounded on $\mathcal{B}_{b_\dagger}$ via the same
$L^\infty$-bounded $(\varphi_\dagger, \sigma_\dagger)$ together with the
inequality $\|e^{-\omega}\|_\infty \le 1/R_\sigma$, which delivers a uniform
$C^2$-boundedness of $\Phi$ on an $L^\infty$ neighbourhood of $b_\dagger$;
the strengthened Assumption~\ref{ass:unifbd}\textup{(iii)} is what supplies
the needed $L^\infty$ control of both components of $b_\dagger$ (alternatively,
under Assumption~\ref{ass:sourcecond} with $\beta > p/2$, Sobolev embedding
delivers the same $L^\infty$ bound).
By the mean-value theorem along the line segment $b_\dagger \to \hat b_{\hat m}$,
\[
\|\hat a_{\hat m} - a_\dagger\|
\;\le\; C_\Phi\,\|\hat b_{\hat m} - b_\dagger\|
\;+\; \tfrac12\,\sup_{b \in \mathcal{B}_{b_\dagger}}\!\|\Phi''(b)\|\,
\|\hat b_{\hat m} - b_\dagger\|^2,
\]
and the quadratic term is bounded explicitly by
\[
\tfrac12\,\sup_{b \in \mathcal B_{b_\dagger}}\!\|\Phi''(b)\|\,
\|\hat b_{\hat m} - b_\dagger\|^{2}
\;\le\;
\tfrac12\,\sup_{b \in \mathcal B_{b_\dagger}}\!\|\Phi''(b)\|
\cdot \mathrm{diam}(\mathcal B_{b_\dagger})
\cdot \|\hat b_{\hat m} - b_\dagger\|,
\]
which is linear in $\|\hat b_{\hat m} - b_\dagger\|$ with a finite constant
depending only on $C^2$-bounds of $\Phi$ on $\mathcal B_{b_\dagger}$ and the
diameter of that ball. The oracle inequality therefore transfers with
constant $C(\tau)\cdot C_\Phi$.

This completes Step~5 and the proof.
\end{proof}

% -------------------------------------------------------------------
% Remarks AFTER the proof (per theorem-environment-style.md).
% -------------------------------------------------------------------

\begin{remark}[Adaptation overhead and the regime distinction]\label{rem:overhead}
% PATCH-M3
The constant $C(\tau)$ in \eqref{eq:lepskiioracle} scales as $5\tau + 1 +
O(\tau^{-1})$, with the choice $\tau = \tau_n = \sqrt{(1+\kappa)\,c_0^{-1}\,\log M_n}$
producing an adaptation overhead of order $\sqrt{\log M_n}$. The overhead is
therefore $\sqrt{\log M_n}$ in all regimes, but specialises as follows. In
the \emph{mildly ill-posed} regime --- $t_j$ polynomial decay, so that
$m^*(n) \asymp n^\alpha$ for some $\alpha \in (0, 1)$ and
$M_n = O(\delta_n^{-1})$ is polynomial in $n$ --- the overhead is
$O(\sqrt{\log n})$. In the \emph{severely ill-posed} regime --- $t_j$
exponential decay, so that $m^*(n) \asymp \log n$ and $M_n = O(\delta_n^{-1})$
is at most polynomial in $\log n$ --- the overhead is
$O(\sqrt{\log\log n})$. The earlier discussion claiming a $\log\log n$
overhead in the severely ill-posed regime was missing the square root;
\citet{hofmannmathe2007} confirm the $\sqrt{\log\log n}$ form in the
exponentially ill-posed case.
\end{remark}

\begin{remark}[Reduction to the linear inverse problem]\label{rem:linear-reduction}
% PATCH-MN7
When the route is MI and $\sigma_\dagger \equiv 1$ is known, the unknown
collapses to $\varphi_\dagger$ alone, the operator equation reduces to the
canonical linear nonparametric IV problem $T\varphi_\dagger = E[Y\mid W]$
with $T\varphi(w) := E[\varphi(X)\mid W = w]$
\citep{darolles2011, hall2005}, the Frechet remainder $III_m$ in the proof
of Lemma~\ref{lem:itervar} vanishes identically, and $K_{b_\dagger,\hat b_j}
= I$ in Proposition~\ref{prp:koper}\textup{(b)} so Term $IV_m$ vanishes.
The variance proxy reduces to
\[
V(m) \;=\; c_v\,\delta_n\sqrt m + c'_v\,\gamma_n\,m^{(1-\beta)/2}\,
\ln(m)^{\un(\beta = 2)}\,m^{(\beta/2-1)\vee 0},
\]
and Proposition~\ref{prp:lepskii}'s oracle inequality recovers two
well-established benchmarks: the noisy-operator Lepskii balancing inequality
of \citet{hoffmannreiss2008} -- with $(\delta_n, \gamma_n)$ corresponding
exactly to data noise on $\hat r := \widehat{E[Y\mid W]}$ and operator noise
on $\hat T$ in their setup -- and the white-noise inverse-problem
adaptation inequality of \citet{cavalier2002} for spectral cutoff. The
genuinely new content of Proposition~\ref{prp:lepskii} beyond these
benchmarks lies in Terms $III_m$ and $IV_m$ of the proof of
Lemma~\ref{lem:itervar}, which arise only when $\sigma_\dagger$ is
unknown and the location-scale structure becomes a genuinely nonlinear
inverse problem.
\end{remark}

\begin{remark}[Extension to $\beta$-mixing data]\label{rem:mixing}
% PATCH-MN8
Assumption \ref{ass:rateconv} is stated for an i.i.d. sample, in line with
the source paper's $\{(Y_i, X_i, W_i)\}_{i=1}^n$. The Bernstein-type
concentration \eqref{eq:bernstein} of Step~1 extends to geometrically
$\beta$-mixing strictly stationary sequences via the kernel-regression
exponential inequality of \citet[Theorem~6]{hansen2008}, which uses
sub-Gaussian conditional response (Assumption~\ref{ass:subgauss}) together
with the standard mixing-rate condition $\beta(k) \le \beta_0\,\rho^k$ for
some $\rho \in (0,1)$. The constants $(c_v, c'_v)$ of
Lemma~\ref{lem:itervar} pick up a finite mixing-dependent factor, but the
form of \eqref{eq:lepskiioracle} is unchanged. This extension covers the
empirical application of Section~\ref{sec:empapp-rates} to T-bill time
series.
\end{remark}

\begin{remark}[Note on the LF spectral-filter constant]\label{rem:cg}
% PATCH-MN3
The LF generating function used implicitly in Step~3 is
$g_m(\lambda) = c_g\sum_{j=0}^{m-1}(1 - c_g\lambda)^j$, where the constant
$c_g$ is the spectral-filter parameter (here $c_g = c$, the LF step size).
The spectral inequality
$\sup_{\lambda \in [0, 1/c_g]}|\lambda^{1/2} g_m(\lambda)| \le \sqrt m$
\citep[Eq.~6.5]{engl2000} uses $c_g$, not the LF step size; we keep the
single symbol $c$ for both quantities throughout the proof for notational
economy, but the absorption $|c \cdot g_m| \le \sqrt m$ behind
Term~$II_m$ should be read with this convention in mind.
\end{remark}

\begin{remark}[Provenance of the propagator substitution]\label{rem:propagator-source}
The propagator lemma (Lemma~\ref{lem:propagator}) makes explicit a step that
is silent in the source paper at the four-term decomposition: it is the
exact analogue, in the operator-valued setting, of the telescoping
identity used in \citet[Section~6.2]{engl2000} and
\citet[Proposition~3.4]{bauerhohage2005} to reduce a Newton-LF iteration
with shifting linearization point to its limit-linearization version plus a
bounded remainder.
\end{remark}
